\documentclass[11pt, a4paper]{article}
\usepackage[utf8]{inputenc}
\usepackage[T1]{fontenc}
\usepackage[spanish]{babel}
\usepackage{amsmath, amssymb}
\usepackage{geometry}
\usepackage{enumitem}
\usepackage{titlesec}
\usepackage{parskip}
\usepackage{mdframed}
\usepackage{xcolor}
\usepackage{array}
\usepackage{booktabs}

\geometry{margin=2.5cm}

\titleformat{\section}{\large\bfseries}{}{0em}{}
\titleformat{\subsection}{\normalsize\bfseries}{}{0em}{}
\titleformat{\subsubsection}{\normalsize\itshape}{}{0em}{}

% Bloque de "respuesta posible"
\newmdenv[
    linecolor=gray!50,
    linewidth=0.5pt,
    backgroundcolor=gray!8,
    skipabove=6pt,
    skipbelow=6pt,
    innertopmargin=6pt,
    innerbottommargin=6pt,
]{respuesta}

\newcommand{\pregunta}[1]{\subsection{#1}}
\newcommand{\eje}[1]{\subsubsection{#1}}
\newcommand{\argmax}{\operatorname*{arg\,m\acute{a}x}}

\title{\textbf{Modelos de Parcial -- Aprendizaje por Refuerzo} \\
\large Inteligencia Artificial y Neurociencias \\ Universidad Torcuato Di Tella}
\author{}
\date{}

\begin{document}
\maketitle

\begin{mdframed}[linecolor=black, linewidth=0.7pt]
\textbf{Formato.} Cada modelo combina \textbf{ejercicios pr\'acticos} (seguimiento manual de algoritmos, c\'alculo de retornos y valores) con \textbf{preguntas de interpretaci\'on} sencillas sobre las f\'ormulas y los conceptos. Los temas cubiertos son: bandidos de $k$ brazos, Procesos de Decisi\'on de Markov, m\'etodos Monte Carlo, m\'etodos de Diferencias Temporales (SARSA y Q-learning), Deep Q-Learning y temas avanzados. Debajo de cada consigna se da una \emph{respuesta posible} a modo de gu\'ia.
\end{mdframed}

%% ============================================================
\section{Modelo 1}
%% ============================================================

\eje{1. Bandidos de $k$ brazos (pr\'actico)}
\pregunta{Se aplica un algoritmo de bandidos de $k=3$ brazos (acciones 1, 2, 3) con selecci\'on $\varepsilon$-greedy, estimando los valores con el promedio muestral y con $Q_1(a)=0$ para toda $a$. La secuencia de acciones y recompensas es: $A_1{=}1, R_1{=}2$; $A_2{=}2, R_2{=}{-}1$; $A_3{=}1, R_3{=}1$; $A_4{=}3, R_4{=}2$; $A_5{=}1, R_5{=}0$. \textquestiondown En qu\'e pasos \textbf{seguro} ocurri\'o el caso $\varepsilon$ (acci\'on al azar) y en qu\'e pasos \textbf{podr\'ia} haber ocurrido?}

\begin{respuesta}
\textit{Respuesta posible.} Se calcula $Q_t(a)$ antes de cada paso y se compara la acci\'on tomada con la(s) acci\'on(es) greedy. Si la acci\'on tomada \emph{no} es greedy, seguro fue exploraci\'on; si lo es, pudo ser greedy o exploraci\'on.

\begin{center}
\begin{tabular}{c c c c l}
\toprule
$t$ & $Q_t(1),Q_t(2),Q_t(3)$ & greedy & $A_t$ & conclusi\'on \\
\midrule
1 & $(0,0,0)$        & $\{1,2,3\}$ & 1 & \textbf{podr\'ia} (empate) \\
2 & $(2,0,0)$        & $\{1\}$     & 2 & \textbf{seguro} $\varepsilon$ \\
3 & $(2,-1,0)$       & $\{1\}$     & 1 & \textbf{podr\'ia} \\
4 & $(1{,}5,-1,0)$   & $\{1\}$     & 3 & \textbf{seguro} $\varepsilon$ \\
5 & $(1{,}5,-1,2)$   & $\{3\}$     & 1 & \textbf{seguro} $\varepsilon$ \\
\bottomrule
\end{tabular}
\end{center}

\noindent (Actualizaciones: tras $t{=}1$, $Q(1){=}2$; tras $t{=}3$, $Q(1){=}\frac{2+1}{2}{=}1{,}5$; tras $t{=}4$, $Q(3){=}2$.)
Entonces: \textbf{seguro} en $t=2,4,5$; \textbf{podr\'ia} en $t=1,3$.
\end{respuesta}

\eje{2. Bandidos (interpretaci\'on)}
\pregunta{La regla de actualizaci\'on incremental puede escribirse como
$\;\text{NuevaEst} = \text{ViejaEst} + \text{TasaApr}\,[\,\text{Objetivo} - \text{ViejaEst}\,]$.
Identifique qu\'e es cada t\'ermino en el caso del promedio muestral ($Q_{n+1}=Q_n+\frac{1}{n}[R_n-Q_n]$) y explique por qu\'e en problemas \textbf{no estacionarios} conviene usar una tasa $\alpha$ constante en vez de $\frac1n$.}

\begin{respuesta}
\textit{Respuesta posible.} \textbf{ViejaEst} es $Q_n$ (la estimaci\'on previa del valor de la acci\'on), \textbf{NuevaEst} es $Q_{n+1}$, el \textbf{Objetivo} es la recompensa reci\'en observada $R_n$, y la \textbf{TasaApr} es $\frac1n$. El t\'ermino $[R_n - Q_n]$ es el \emph{error de estimaci\'on}: la actualizaci\'on mueve la estimaci\'on una fracci\'on de ese error hacia el objetivo.

Con $\frac1n$, la tasa decrece con el tiempo y todas las recompensas pesan igual (promedio muestral), lo que converge al valor real \emph{si} el valor no cambia. Pero si el problema es \textbf{no estacionario} (los valores $q_*(a)$ var\'ian en el tiempo), interesa dar \textbf{m\'as peso a las recompensas recientes}. Con $\alpha$ constante, $Q_{n+1}=(1-\alpha)^n Q_1 + \sum_{i=1}^{n}\alpha(1-\alpha)^{n-i}R_i$: los pesos de las recompensas viejas decaen exponencialmente, as\'i la estimaci\'on ``olvida'' el pasado y sigue los cambios del ambiente.
\end{respuesta}

\eje{3. MDP: retornos (pr\'actico)}
\pregunta{Dada la secuencia de recompensas $R_1{=}2,\,R_2{=}0,\,R_3{=}{-}1,\,R_4{=}3,\,R_5{=}1$ con $T=5$ y $\gamma=0{,}5$, calcular $G_0, G_1, \dots, G_4$. \emph{Ayuda: de atr\'as hacia adelante usando $G_t = R_{t+1} + \gamma\,G_{t+1}$.}}

\begin{respuesta}
\textit{Respuesta posible.} Con $G_5 = 0$:
\begin{align*}
G_4 &= R_5 + \gamma G_5 = 1 + 0{,}5\cdot 0 = 1 \\
G_3 &= R_4 + \gamma G_4 = 3 + 0{,}5\cdot 1 = 3{,}5 \\
G_2 &= R_3 + \gamma G_3 = -1 + 0{,}5\cdot 3{,}5 = 0{,}75 \\
G_1 &= R_2 + \gamma G_2 = 0 + 0{,}5\cdot 0{,}75 = 0{,}375 \\
G_0 &= R_1 + \gamma G_1 = 2 + 0{,}5\cdot 0{,}375 = 2{,}1875
\end{align*}
\end{respuesta}

\eje{4. MDP (interpretaci\'on)}
\pregunta{Explique qu\'e representa el \textbf{factor de descuento} $\gamma$ en $G_t = \sum_{k=t+1}^{T}\gamma^{\,k-t-1}R_k$. \textquestiondown Qu\'e comportamiento induce $\gamma=0$ y qu\'e comportamiento induce $\gamma\to 1$? Enuncie adem\'as la \textbf{propiedad de Markov}.}

\begin{respuesta}
\textit{Respuesta posible.} $\gamma \in [0,1]$ determina cu\'anto valora el agente las recompensas futuras frente a las inmediatas. Con \textbf{$\gamma=0$} el agente es \emph{miope}: solo le importa la recompensa inmediata $R_{t+1}$. Con \textbf{$\gamma\to 1$} valora casi por igual recompensas lejanas y cercanas (agente ``previsor''). Sirve la analog\'ia con la inflaci\'on: una inflaci\'on alta (que licua el valor del dinero futuro) equivale a un $\gamma$ bajo.

La \textbf{propiedad de Markov} dice que la probabilidad del pr\'oximo estado y recompensa depende \emph{solo} del estado y la acci\'on actuales, y no de la historia previa: $p(s',r\mid s,a)=\Pr\{S_t{=}s',R_t{=}r \mid S_{t-1}{=}s, A_{t-1}{=}a\}$. Esto obliga a definir los estados de modo que incluyan toda la informaci\'on del pasado relevante para el futuro.
\end{respuesta}

\eje{5. Q-learning (pr\'actico)}
\pregunta{Ambiente con estados $\{s_1, s_2\}$ y acciones $\{\leftarrow, \rightarrow\}$ (tarea continua). Se ejecuta Q-learning con $\alpha=0{,}5$, $\gamma=0{,}5$ y $Q$ inicializada en 0. La trayectoria es: $s_1, \rightarrow, +4, s_2, \rightarrow, +2, s_1, \rightarrow, +8, s_2, \dots$ Hacer el seguimiento de $Q$ tras cada acci\'on y dar la pol\'itica greedy resultante.}

\begin{respuesta}
\textit{Respuesta posible.} Regla: $Q(S,A)\leftarrow Q(S,A) + \alpha\big[R + \gamma\max_a Q(S',a) - Q(S,A)\big]$.

\textbf{Paso 1} ($s_1,\rightarrow,+4,s_2$): $\max_a Q(s_2,a)=0$.
$\;Q(s_1,\rightarrow)=0+0{,}5[\,4+0{,}5\cdot 0-0\,]=\mathbf{2}$.

\textbf{Paso 2} ($s_2,\rightarrow,+2,s_1$): $\max_a Q(s_1,a)=\max(2,0)=2$.
$\;Q(s_2,\rightarrow)=0+0{,}5[\,2+0{,}5\cdot 2-0\,]=\mathbf{1{,}5}$.

\textbf{Paso 3} ($s_1,\rightarrow,+8,s_2$): $\max_a Q(s_2,a)=\max(1{,}5,0)=1{,}5$.
$\;Q(s_1,\rightarrow)=2+0{,}5[\,8+0{,}5\cdot 1{,}5-2\,]=2+0{,}5\cdot 6{,}75=\mathbf{5{,}375}$.

Tabla final: $Q(s_1,\rightarrow)=5{,}375$, $Q(s_2,\rightarrow)=1{,}5$, resto $=0$.
\textbf{Pol\'itica greedy}: en $s_1$ elegir $\rightarrow$; en $s_2$ elegir $\rightarrow$.
\end{respuesta}

\eje{6. TD (interpretaci\'on)}
\pregunta{Escriba las reglas de actualizaci\'on de \textbf{SARSA} y \textbf{Q-learning} y explique cu\'al es la \'unica diferencia entre ambas. \textquestiondown Por qu\'e se dice que SARSA es \emph{on-policy} y Q-learning \emph{off-policy}?}

\begin{respuesta}
\textit{Respuesta posible.}
\[
\text{SARSA:}\quad Q(S,A)\leftarrow Q(S,A)+\alpha\big[R+\gamma\,Q(S',A')-Q(S,A)\big]
\]
\[
\text{Q-learning:}\quad Q(S,A)\leftarrow Q(S,A)+\alpha\big[R+\gamma\,\max_a Q(S',a)-Q(S,A)\big]
\]
La \'unica diferencia est\'a en c\'omo estiman el retorno desde $S'$: SARSA usa $Q(S',A')$, donde $A'$ es la acci\'on que \textbf{realmente} va a tomar seg\'un su pol\'itica $\varepsilon$-greedy; Q-learning usa $\max_a Q(S',a)$, es decir asume una pol\'itica \textbf{greedy pura} para estimar el retorno.

Por eso SARSA es \textbf{on-policy}: eval\'ua y mejora la misma pol\'itica que usa para actuar (incluye el efecto de la exploraci\'on). Q-learning es \textbf{off-policy}: la pol\'itica que estima (greedy) est\'a \emph{desacoplada} de la que sigue para explorar ($\varepsilon$-greedy).
\end{respuesta}

%% ============================================================
\section{Modelo 2}
%% ============================================================

\eje{1. Bandidos: actualizaci\'on incremental (pr\'actico)}
\pregunta{Un algoritmo de bandidos con promedio muestral tiene, para la acci\'on $a$, la estimaci\'on $Q_n(a)=3$ luego de haberla tomado $n-1=4$ veces. Se vuelve a elegir $a$ y se obtiene $R_n = 8$. (a)~Calcular $Q_{n+1}(a)$ con la f\'ormula incremental. (b)~\textquestiondown Por qu\'e esta forma es preferible a guardar todas las recompensas y promediar?}

\begin{respuesta}
\textit{Respuesta posible.} (a) $n=5$, entonces
\[
Q_{n+1}(a) = Q_n + \tfrac{1}{n}\big[R_n - Q_n\big] = 3 + \tfrac{1}{5}\,[\,8-3\,] = 3 + 1 = 4.
\]
(b) La actualizaci\'on incremental requiere almacenar solo $Q_n$ y $n$ (memoria $O(1)$) y cada actualizaci\'on cuesta $O(1)$, en vez de guardar la lista completa de recompensas y recomputar la media (memoria y tiempo crecientes con la cantidad de muestras).
\end{respuesta}

\eje{2. Bandidos (interpretaci\'on)}
\pregunta{\textquestiondown Es equivalente usar una pol\'itica $\varepsilon$-greedy con $\varepsilon=0{,}1$ durante 1000 pasos a primero ejecutar 100 pasos puramente exploratorios y luego 900 pasos puramente greedy? Justifique. Explique tambi\'en la idea de los \textbf{valores iniciales optimistas} como forma de exploraci\'on.}

\begin{respuesta}
\textit{Respuesta posible.} \textbf{No} son equivalentes. Con $\varepsilon$-greedy la exploraci\'on est\'a \textbf{distribuida} a lo largo de todo el proceso: sigue explorando (aprox.\ un 10\,\% del tiempo) incluso al final, lo que le permite corregir estimaciones err\'oneas y adaptarse (\'util si el problema no es estacionario). El esquema ``explorar primero, explotar despu\'es'' fija las estimaciones tras la fase inicial: si por azar una buena acci\'on qued\'o subestimada en esos primeros pasos, ya nunca la vuelve a probar. Adem\'as, durante la fase greedy no obtiene m\'as informaci\'on de las acciones no elegidas.

\textbf{Valores iniciales optimistas}: se inicializa $Q_1(a)$ con un valor muy alto para toda $a$. Como toda recompensa real resulta ``decepcionante'' frente a la estimaci\'on inflada, el error $[R-Q]$ es negativo y el agente se ve \emph{empujado} a probar las otras acciones (que siguen con valor alto). As\'i se logra una exploraci\'on inicial sistem\'atica sin necesidad de aleatoriedad. Su limitaci\'on: es un mecanismo transitorio, poco \'util en problemas no estacionarios.
\end{respuesta}

\eje{3. MDP: funciones de valor (pr\'actico)}
\pregunta{Ambiente determin\'istico en l\'inea: $A \to B \to C \to D$, con $D$ terminal. Recompensas: $A\to B$ da $1$, $B\to C$ da $2$, $C\to D$ da $3$. Con la pol\'itica $\pi$ que avanza siempre a la derecha, calcular $v_\pi(A), v_\pi(B), v_\pi(C)$ para $\gamma = 0$, $\gamma = 0{,}5$ y $\gamma = 1$.}

\begin{respuesta}
\textit{Respuesta posible.} Usando $v_\pi(s) = r + \gamma\, v_\pi(s')$ desde el final ($v_\pi(D)=0$):
\[
v_\pi(C)=3+\gamma\cdot 0=3,\qquad v_\pi(B)=2+\gamma\, v_\pi(C),\qquad v_\pi(A)=1+\gamma\, v_\pi(B).
\]
\begin{center}
\begin{tabular}{c c c c}
\toprule
$\gamma$ & $v_\pi(A)$ & $v_\pi(B)$ & $v_\pi(C)$ \\
\midrule
$0$   & $1$    & $2$   & $3$ \\
$0{,}5$ & $2{,}75$ & $3{,}5$ & $3$ \\
$1$   & $6$    & $5$   & $3$ \\
\bottomrule
\end{tabular}
\end{center}
(Con $\gamma=0{,}5$: $v_\pi(B)=2+0{,}5\cdot3=3{,}5$; $v_\pi(A)=1+0{,}5\cdot3{,}5=2{,}75$.
Con $\gamma=1$: $v_\pi(B)=5$, $v_\pi(A)=6$.) Como $\pi$ es determin\'istica, $q_\pi(s,\rightarrow)=v_\pi(s)$ para la acci\'on seguida.
\end{respuesta}

\eje{4. Monte Carlo (pr\'actico)}
\pregunta{Ambiente con estados $\{s_1, s_2, s_3\}$ ($s_3$ terminal) y acciones $\{a,b\}$. Se corre Monte Carlo para estimar $\pi_*$ con $\gamma=1$ y $Q$ inicializada en 0. Episodios:
\emph{Ep.\ 1}: $s_1, a, {+}2, s_2, a, {+}3, s_3$; \quad \emph{Ep.\ 2}: $s_1, a, {+}1, s_2, b, {+}4, s_3$.
Indicar $\mathrm{Retornos}(S,A)$ y $Q(S,A)$ al final de cada episodio, y la acci\'on greedy en cada estado al comenzar el tercer episodio.}

\begin{respuesta}
\textit{Respuesta posible.} Los retornos ($\gamma=1$) son la suma de recompensas desde cada par $(S,A)$.

\emph{Fin Ep.\ 1}: $G(s_1,a)=2+3=5$, $G(s_2,a)=3$.
$\;\mathrm{Ret}(s_1,a)=[5]\Rightarrow Q(s_1,a)=5$;\; $\mathrm{Ret}(s_2,a)=[3]\Rightarrow Q(s_2,a)=3$.

\emph{Fin Ep.\ 2}: $G(s_1,a)=1+4=5$, $G(s_2,b)=4$.
$\;\mathrm{Ret}(s_1,a)=[5,5]\Rightarrow Q(s_1,a)=5$;\; $\mathrm{Ret}(s_2,b)=[4]\Rightarrow Q(s_2,b)=4$.

Tabla: $Q(s_1,a)=5,\ Q(s_1,b)=0,\ Q(s_2,a)=3,\ Q(s_2,b)=4$.
\textbf{Greedy} al comenzar Ep.\ 3: en $s_1$ elegir $a$ ($5>0$); en $s_2$ elegir $b$ ($4>3$).
\end{respuesta}

\eje{5. MC vs TD (interpretaci\'on)}
\pregunta{\textquestiondown Cu\'al es la diferencia fundamental entre los m\'etodos Monte Carlo y los de Diferencias Temporales en cu\'ando y c\'omo actualizan $Q$? Defina \textbf{bootstrapping} y d\'e una situaci\'on donde TD tenga ventaja pr\'actica sobre MC.}

\begin{respuesta}
\textit{Respuesta posible.} \textbf{Monte Carlo} espera al \textbf{final del episodio} y actualiza usando el retorno real medido $G_t$: $Q(S_t,A_t)\leftarrow Q(S_t,A_t)+\alpha[G_t - Q(S_t,A_t)]$. \textbf{TD} actualiza \textbf{tras cada paso}, sin esperar el final, usando $R_{t+1}+\gamma Q(S_{t+1},A_{t+1})$ como objetivo.

\textbf{Bootstrapping}: estimar un valor a partir de \emph{otras estimaciones} propias (aqu\'i, $Q(S_{t+1},\cdot)$) en lugar de esperar el resultado real completo. TD hace bootstrapping; MC no.

\textbf{Ventaja de TD}: en tareas \textbf{continuas} (sin fin) o con episodios muy largos, MC no podr\'ia actualizar (no hay retorno final que medir), mientras que TD aprende ``sobre la marcha''. Ejemplo: un agente en un videojuego con muchos estados aprende en tiempo real con TD, mientras que MC tendr\'ia que completar la partida entera para cada actualizaci\'on.
\end{respuesta}

\eje{6. Deep Q-Learning (interpretaci\'on)}
\pregunta{\textquestiondown Por qu\'e los m\'etodos tabulares no sirven en problemas como jugar al Atari desde los p\'ixeles? \textquestiondown Qu\'e se hace en su lugar? Explique brevemente para qu\'e sirven el \textbf{replay de experiencias} y la \textbf{target network} en DQN.}

\begin{respuesta}
\textit{Respuesta posible.} Los m\'etodos tabulares guardan un valor por cada par $(s,a)$ en una tabla $Q$, lo cual es imposible cuando el espacio de estados $\mathcal{S}$ es gigantesco: la cantidad de pantallas posibles de un juego es combinatoriamente enorme, y adem\'as dos pantallas casi id\'enticas cuentan como estados distintos (no hay generalizaci\'on). En su lugar se usa \textbf{aproximaci\'on de la funci\'on de valor}: se reemplaza la tabla por una funci\'on parametrizada $\hat q(s,a,\mathbf{w})$ (una red neuronal con pesos $\mathbf{w}$, $d \ll |\mathcal{S}|$) que se ajusta con descenso por gradiente (semi-gradiente), permitiendo generalizar a estados nunca vistos.

\textbf{Replay de experiencias}: se almacenan las transiciones $e_t=(S_t,A_t,R_{t+1},S_{t+1})$ en una memoria y cada actualizaci\'on se hace sobre un \emph{minibatch} muestreado al azar; esto rompe la correlaci\'on temporal entre muestras consecutivas y estabiliza el entrenamiento. \textbf{Target network}: una copia de la red que se usa para computar el objetivo y se actualiza solo cada cierta cantidad de iteraciones; al mantener el objetivo ``fijo'' un rato, evita que la red persiga un blanco que se mueve con cada paso.
\end{respuesta}

%% ============================================================
\section{Modelo 3}
%% ============================================================

\eje{1. Bandidos: UCB (pr\'actico)}
\pregunta{Se usa selecci\'on por l\'imite de confianza superior (UCB): $A_t \doteq \argmax_a\big[Q_t(a) + c\sqrt{\ln t / N_t(a)}\big]$, con $c=1$. Ya se ejecutaron 4 pasos ($t=5$, $\ln 5 \approx 1{,}609$) y el estado es $N(1){=}2, Q(1){=}1$; $N(2){=}1, Q(2){=}2$; $N(3){=}1, Q(3){=}0$. \textquestiondown Qu\'e acci\'on elige UCB en $t=5$?}

\begin{respuesta}
\textit{Respuesta posible.}
\begin{align*}
\text{UCB}(1) &= 1 + \sqrt{1{,}609/2} = 1 + 0{,}897 \approx 1{,}90 \\
\text{UCB}(2) &= 2 + \sqrt{1{,}609/1} = 2 + 1{,}269 \approx 3{,}27 \\
\text{UCB}(3) &= 0 + \sqrt{1{,}609/1} = 1{,}269 \approx 1{,}27
\end{align*}
El m\'aximo es la \textbf{acci\'on 2}. (El t\'ermino $c\sqrt{\ln t/N_t(a)}$ crece cuando una acci\'on se explor\'o poco: premia la \emph{incertidumbre}, favoreciendo acciones con pocas visitas.)
\end{respuesta}

\eje{2. MDP: retorno infinito (pr\'actico)}
\pregunta{(a) Con $\gamma=0{,}8$ y la secuencia $R_1=3$ seguida de una sucesi\'on infinita de dieces ($R_2=R_3=\dots=10$), calcular $G_0$. (b) Demostrar que si $R_t=1$ para todo $t$, entonces $G_t = \frac{1}{1-\gamma}$.}

\begin{respuesta}
\textit{Respuesta posible.} (a) Desde $t=1$, la cola es una serie geom\'etrica: $G_1 = \sum_{k=0}^{\infty}10\,\gamma^{k} = \frac{10}{1-0{,}8}=50$. Entonces $G_0 = R_1 + \gamma G_1 = 3 + 0{,}8\cdot 50 = \mathbf{43}$.

(b) $G_t = \sum_{k=0}^{\infty}\gamma^{k}\cdot 1 = 1 + \gamma + \gamma^2 + \dots = \frac{1}{1-\gamma}$ (serie geom\'etrica de raz\'on $\gamma<1$). Alternativamente, por la forma recursiva $G_t = 1 + \gamma G_t \Rightarrow G_t(1-\gamma)=1 \Rightarrow G_t=\frac{1}{1-\gamma}$.
\end{respuesta}

\eje{3. SARSA vs Q-learning en una transici\'on (pr\'actico)}
\pregunta{La tabla $Q$ tiene: $Q(0,a)=4$; en el estado $1$: $Q(1,x)=6, Q(1,y)={-}2, Q(1,z)=1$. El agente est\'a en $S=0$, ejecuta $a$, recibe $R=3$ y pasa a $S'=1$. Con $\gamma=0{,}5$ y $\alpha=0{,}5$: (a)~actualizar $Q(0,a)$ por \textbf{Q-learning}; (b)~actualizar $Q(0,a)$ por \textbf{SARSA} considerando todas las opciones de $A'$ (explotaci\'on y cada exploraci\'on posible).}

\begin{respuesta}
\textit{Respuesta posible.} (a) \textbf{Q-learning}: usa $\max_a Q(1,a)=6$.
\[
Q(0,a) \leftarrow 4 + 0{,}5\,[\,3 + 0{,}5\cdot 6 - 4\,] = 4 + 0{,}5\cdot 2 = \mathbf{5}.
\]
(b) \textbf{SARSA}: depende de la acci\'on $A'$ que se elija en $S'=1$.
\begin{itemize}[nosep]
\item Explotaci\'on / exploraci\'on que cae en $x$ ($Q=6$): $4+0{,}5[3+0{,}5\cdot6-4]=\mathbf{5}$.
\item Exploraci\'on en $y$ ($Q=-2$): $4+0{,}5[3+0{,}5\cdot(-2)-4]=4+0{,}5\cdot(-2)=\mathbf{3}$.
\item Exploraci\'on en $z$ ($Q=1$): $4+0{,}5[3+0{,}5\cdot1-4]=4+0{,}5\cdot(-0{,}5)=\mathbf{3{,}75}$.
\end{itemize}
Q-learning siempre toma el m\'aximo; SARSA usa el valor de la acci\'on realmente elegida, por eso su resultado var\'ia con la exploraci\'on.
\end{respuesta}

\eje{4. Q-learning: par\'ametros (interpretaci\'on)}
\pregunta{En Q-learning, \textquestiondown qu\'e efecto tiene usar valores extremos en cada par\'ametro? Analice: (a) $\alpha=0$ y $\alpha=1$; (b) $\gamma=0$ y $\gamma=1$; (c) $\varepsilon=0$ y $\varepsilon=1$.}

\begin{respuesta}
\textit{Respuesta posible.}
\textbf{(a) $\alpha$ (tasa de aprendizaje).} Con $\alpha=0$ la actualizaci\'on es nula: $Q$ nunca cambia, no aprende. Con $\alpha=1$ la nueva estimaci\'on reemplaza por completo a la vieja ($Q(S,A)=R+\gamma\max_a Q(S',a)$): aprende solo de la \'ultima experiencia, \'util en ambientes determin\'isticos pero inestable si hay ruido.

\textbf{(b) $\gamma$ (descuento).} Con $\gamma=0$ el agente es miope: solo maximiza la recompensa inmediata, ignora el futuro. Con $\gamma=1$ pesa todas las recompensas futuras por igual; en tareas continuas el retorno puede diverger.

\textbf{(c) $\varepsilon$ (exploraci\'on).} Con $\varepsilon=0$ es puramente greedy: nunca explora, puede quedar atrapado en una pol\'itica sub\'optima. Con $\varepsilon=1$ act\'ua siempre al azar: explora todo pero nunca aprovecha lo aprendido (aunque, al ser off-policy, Q-learning igual puede estimar la pol\'itica \'optima si visita suficientemente todos los pares).
\end{respuesta}

\eje{5. Ambiente estoc\'astico vs determin\'istico (interpretaci\'on)}
\pregunta{En el juego ``Cuatro en l\'inea'' (dos jugadores), \textquestiondown la funci\'on $p(s',r\mid s,a)$ que caracteriza el ambiente es estoc\'astica o determin\'istica? Justifique. \textquestiondown Qu\'e significa que $\sum_{s',r} p(s',r\mid s,a)=1$?}

\begin{respuesta}
\textit{Respuesta posible.} Desde la perspectiva del agente, el ambiente incluye al \textbf{oponente}. Si el oponente no juega siempre lo mismo (tiene una pol\'itica estoc\'astica o desconocida), entonces tras la acci\'on del agente el estado resultante $s'$ (que depende de la respuesta del rival) no es \'unico: el ambiente es \textbf{estoc\'astico}. Si el oponente fuera totalmente determin\'istico y conocido, $p$ ser\'ia determin\'istica. La regla del juego en s\'i es determin\'istica, pero la din\'amica $p$ vista por el agente hereda la incertidumbre del rival.

$\sum_{s',r} p(s',r\mid s,a)=1$ expresa que, fijados un estado $s$ y una acci\'on $a$, la distribuci\'on sobre todos los posibles pares (estado siguiente, recompensa) es una distribuci\'on de probabilidad v\'alida: \emph{algo} ocurre con certeza, y las probabilidades de todos los resultados posibles suman 1.
\end{respuesta}

\eje{6. Temas avanzados (interpretaci\'on)}
\pregunta{Los m\'etodos de \textbf{gradiente de pol\'iticas} (p.ej.\ REINFORCE) se contraponen a los basados en funci\'on de valor. \textquestiondown Qu\'e aprenden directamente y en qu\'e casos conviene usarlos? Explique adem\'as, en una l\'inea, qu\'e hace \textbf{Monte Carlo Tree Search (MCTS)}.}

\begin{respuesta}
\textit{Respuesta posible.} Los m\'etodos de gradiente de pol\'iticas aprenden \textbf{directamente una pol\'itica parametrizada} $\pi(a\mid s,\boldsymbol\theta)$ (las probabilidades de cada acci\'on), sin pasar por una funci\'on de valor $q_*$. Se ajustan los par\'ametros por ascenso de gradiente sobre una m\'etrica de desempe\~no $J(\boldsymbol\theta)$. En REINFORCE, la probabilidad de una acci\'on \textbf{sube o baja proporcionalmente al retorno observado} $G$. Convienen cuando hay \textbf{demasiadas acciones}, acciones \textbf{continuas}, o cuando es m\'as simple representar la pol\'itica que el valor (operaci\'on de brazos rob\'oticos, trading, LLMs).

\textbf{MCTS}: planificaci\'on en tiempo de decisi\'on que construye incrementalmente un \'arbol desde el estado actual, simulando trayectorias (selecci\'on--expansi\'on--simulaci\'on--backup) para estimar el valor de las acciones y elegir la mejor. Es el motor de b\'usqueda de AlphaGo.
\end{respuesta}

%% ============================================================
\section{Banco de preguntas extra (para practicar)}
%% ============================================================

\subsection*{Bandidos de $k$ brazos}
\begin{itemize}[nosep]
    \item Explique el compromiso (\emph{trade-off}) entre \textbf{exploraci\'on} y \textbf{explotaci\'on}. \textquestiondown Qu\'e es una acci\'on \emph{greedy}?
    \item Diferencia entre $q_*(a)$ (valor real) y $Q_t(a)$ (estimaci\'on). \textquestiondown Por qu\'e $Q_t(a)\to q_*(a)$ con el promedio muestral? (Ley de los Grandes N\'umeros).
    \item En el ejemplo de 10 brazos, \textquestiondown por qu\'e $\varepsilon=0{,}1$ termina eligiendo m\'as la acci\'on \'optima que $\varepsilon=0$ (greedy puro)?
\end{itemize}

\subsection*{Procesos de Decisi\'on de Markov}
\begin{itemize}[nosep]
    \item Distinga $v_\pi(s)$ de $q_\pi(s,a)$ y escriba la relaci\'on $q_\pi(s,a)=\mathbb{E}_\pi[R_{t+1}+\gamma v_\pi(S_{t+1})\mid S_t{=}s, A_t{=}a]$.
    \item Exprese $p(s'\mid s,a)$ y $r(s,a)$ en funci\'on de $p(s',r\mid s,a)$ (marginalizando y tomando esperanza).
    \item Diferencia entre m\'etodos \textbf{model-based} (p.ej.\ Value Iteration) y \textbf{model-free} (Monte Carlo, Q-learning). \textquestiondown Por qu\'e casi siempre usamos model-free?
    \item \textquestiondown Qu\'e es una pol\'itica \'optima $\pi_*$ y las funciones \'optimas $v_*$, $q_*$?
    \item Modele el Tatet\'i como MDP: estados, acciones, recompensas, \textquestiondown es determin\'istico o estoc\'astico?
\end{itemize}

\subsection*{Monte Carlo y Diferencias Temporales}
\begin{itemize}[nosep]
    \item \textquestiondown Por qu\'e el primer algoritmo MC para estimar $\pi_*$ necesita una pol\'itica $\varepsilon$-greedy y no una greedy pura? (Para visitar todos los pares $(s,a)$).
    \item Escriba la actualizaci\'on incremental de la media de retornos: $Q \leftarrow Q + \frac{1}{k+1}(G^{(k+1)}-Q)$. \textquestiondown Con qu\'e regla de bandidos se parece?
    \item Efecto de arrancar con $\varepsilon$ alto (p.ej.\ 1) y bajarlo gradualmente hasta $\approx 0{,}01$ en MC. (Explorar mucho al principio, explotar al final).
    \item En el ejemplo ``conduciendo a casa'', \textquestiondown por qu\'e TD corrige la predicci\'on en cada tramo y MC reci\'en al llegar?
\end{itemize}

\subsection*{Deep Q-Learning y temas avanzados}
\begin{itemize}[nosep]
    \item \textquestiondown Qu\'e es el error cuadr\'atico $\overline{\mathrm{QE}}(\mathbf{w})=\sum_{s,a}[q_*(s,a)-\hat q(s,a,\mathbf{w})]^2$ y c\'omo se lo reduce con SGD?
    \item \textquestiondown Por qu\'e DQL se llama m\'etodo \emph{semi-gradiente} y por qu\'e no garantiza convergencia (a diferencia del caso tabular)?
    \item M\'etodos \textbf{Actor-Critic}: \textquestiondown qu\'e rol cumple el \emph{actor} y qu\'e rol el \emph{cr\'itico}? \textquestiondown En qu\'e se diferencia de REINFORCE? (Bootstrapping / objetivo = error TD).
    \item Los cinco pasos de MCTS: selecci\'on, expansi\'on, simulaci\'on, backup, repetir. Diferencia entre \emph{pol\'itica del \'arbol} y \emph{pol\'itica rollout}.
    \item \textbf{AlphaGo} vs \textbf{AlphaGo Zero}: \textquestiondown en qu\'e se diferencian respecto del uso de datos humanos y del \emph{self-play} (\emph{tabula rasa})?
    \item Preguntas para neurociencias: \textquestiondown tiene el cerebro un mecanismo parecido a $\varepsilon$-greedy? \textquestiondown C\'omo se relaciona la recompensa de RL con la dopamina? \textquestiondown Hay algo an\'alogo al \emph{replay de experiencias} durante el sue\~no?
\end{itemize}

\end{document}
